\documentclass[12pt,a4paper]{article}
\usepackage[utf8]{inputenc}
\usepackage{ctex}
\usepackage{fontspec}
\usepackage{geometry}
\usepackage{graphicx}
\usepackage{amsmath}
\usepackage{amssymb}
\usepackage{booktabs}
\usepackage{array}
\usepackage{float}
\usepackage{hyperref}
\usepackage{enumerate}

\geometry{left=2.5cm,right=2.5cm,top=2.5cm,bottom=2.5cm}

\title{\LARGE\bfseries 碳化硅外延层厚度的确定}
\author{数学建模竞赛参赛队}
\date{\today}

\begin{document}
\maketitle

\begin{abstract}
本研究针对碳化硅(SiC)外延层厚度测量问题，提出了一种基于红外干涉光谱法的精确厚度计算方法。通过分析SiC外延层/SiC衬底及Si外延层/Si衬底结构在红外波段的干涉光谱，利用双光束干涉原理建立厚度计算模型，其中SiC层采用4H-SiC Sellmeier色散公式计算波长依赖的折射率，Si层采用常数折射率$n=3.42$。对附件提供的4个样品的干涉光谱数据进行峰值检测与统计分析，得到各样品外延层厚度为：SiC-1样品$24.58\,\mu\text{m}$，SiC-2样品$24.02\,\mu\text{m}$，Si-1样品$3.67\,\mu\text{m}$，Si-2样品$3.52\,\mu\text{m}$。实验结果表明，该方法具有较高的测量精度，可为SiC和Si外延片的质量检测提供可靠依据。

\textbf{关键词}: 红外干涉；外延层厚度；干涉峰间距；Sellmeier色散公式；峰值检测
\end{abstract}

\section{引言}

碳化硅(SiC)作为第三代半导体材料的代表，具有宽带隙、高热导率、高击穿电场等优异特性，在电力电子器件领域应用前景广阔。外延层厚度是影响器件性能的关键参数之一，精确测量SiC外延层厚度对于工艺控制和器件设计具有重要意义。

传统的厚度测量方法包括台阶仪测量、光学显微镜测量等，但这些方法存在测量速度慢、对样品有损或精度有限等问题。红外干涉法是一种非接触、无损的高精度测量方法，其基本原理是利用干涉光谱中相邻干涉峰之间的波数间距反推薄膜厚度。当光波垂直入射时，相邻干涉极大值对应的波数间距$\Delta\sigma$与薄膜厚度$d$满足关系：

\begin{equation}
d = \frac{1}{2 \cdot n \cdot \Delta\sigma}
\end{equation}

其中$n$为薄膜材料的折射率，$\Delta\sigma$为相邻干涉峰的波数间距(单位：cm$^{-1}$)，$d$的单位为cm。

\section{数据探索与预处理}

附件提供了4个样品的红外反射光谱数据：
\begin{itemize}
\item 附件1、附件2：SiC外延层/SiC衬底结构，SiC外延层厚度待测
\item 附件3、附件4：Si外延层/Si衬底结构，Si外延层厚度待测
\end{itemize}

数据格式为两列：波数(单位cm$^{-1}$)和反射率(单位$\%$)。首先对原始数据进行探索性分析，观察干涉峰的分布特征，确定有效的分析波段范围。

\subsection{SiC样品分析波段确定}
4H-SiC在约$830\,\text{cm}^{-1}$处存在晶格吸收峰，该波段附近折射率变化剧烈，不宜作为分析区域。经文献调研，$900\sim1300\,\text{cm}^{-1}$波段SiC材料透明性好，折射率稳定，适合作为分析区域。

\subsection{Si样品分析波段确定}
Si在红外波段(波数$400\sim3500\,\text{cm}^{-1}$)透明度良好，折射率近似为常数，可采用宽波段分析以获得更多干涉级次，提高测量精度。

\section{模型假设与符号说明}

\subsection{模型假设}
\begin{enumerate}
\item 光波近似垂直入射(入射角$\theta \approx 0^{\circ}$)，折射角$\theta_t \approx 0^{\circ}$，公式简化为$d = 1/(2n\Delta\sigma)$。
\item SiC层折射率采用4H-SiC Sellmeier色散公式计算：
\begin{equation}
n(\lambda) = \sqrt{6.7\left(1+0.46\frac{\lambda^2}{\lambda^2-0.106^2}\right)}
\end{equation}
其中$\lambda$为真空中的波长(单位：$\mu\text{m}$)，可通过$\lambda = 10^4/\sigma$由波数$\sigma$计算得到。
\item Si层折射率取常数$n_{\text{Si}} = 3.42$(红外波段近似)。
\item 干涉光谱信噪比足够高，干涉峰清晰可辨。
\end{enumerate}

\subsection{符号说明}
\begin{table}[H]
\centering
\begin{tabular}{ccc}
\toprule
\textbf{符号} & \textbf{意义} & \textbf{说明} \
\midrule
$d$ & 外延层厚度 & 单位：cm \
$\Delta\sigma$ & 相邻干涉峰波数间距 & 单位：cm$^{-1}$ \
$n$ & 折射率 & 与波长相关 \
$\sigma$ & 波数 & 单位：cm$^{-1}$ \
$\lambda$ & 波长 & 单位：$\mu\text{m}$ \
$K$ & 干涉对比度 & $K>0.5$时为多光束干涉 \
\bottomrule
\end{tabular}
\caption{主要符号说明}
\end{table}

\section{模型建立}

\subsection{双光束干涉模型}
对于典型的双光束干涉(Newton环、Fabry-Perot标准具)，干涉条件为：
\begin{equation}
2 n d \cos\theta_t = m \lambda \quad (m = 1,2,3,\dots)
\end{equation}

由波数$\sigma=1/\lambda$可得相邻干涉级次间的波数间距：
\begin{equation}
\Delta\sigma = \frac{1}{2 n d \cos\theta_t}
\end{equation}

当$\theta_t \approx 0$时，$\cos\theta_t \approx 1$，得简化公式：
\begin{equation}
d = \frac{1}{2 n \Delta\sigma}
\end{equation}

\subsection{多光束Fabry-Perot干涉模型}
当界面反射率较高时(如Si/空气界面$R\approx 30\%$)，形成多光束干涉，干涉峰更锐利(半高宽更小)，但$\Delta\sigma$与$d$的关系形式仍为上式，区别在于干涉级次$m$更高、峰数更多。

\subsection{折射率计算}
4H-SiC为六方纤锌矿结构，其Sellmeier色散公式为：
\begin{equation}
n^2(\lambda) = 6.7\left(1 + 0.46\frac{\lambda^2}{\lambda^2 - 0.106^2}\right)
\end{equation}

在$900\sim1300\,\text{cm}^{-1}$波段，取$\sigma \approx 1100\,\text{cm}^{-1}$，则$\lambda = 10^4/1100 \approx 9.09\,\mu\text{m}$，代入得$n_{\text{SiC}} \approx 2.68$。
